% Den förenklade kontrollerns block plus lagret som gör den nåbar.
\begin{tikzpicture}[
  every node/.style={font=\scriptsize},
  blk/.style={draw=rule, rounded corners=2pt, minimum width=2.6cm, minimum height=0.72cm,
              align=center, font=\scriptsize, fill=white},
  grp/.style={draw=rule, dashed, rounded corners=3pt, inner sep=6pt},
  arr/.style={-{Stealth[length=1.8mm]}, muted, line width=0.5pt}]

  % Mjukvarusidan.
  \node[blk, fill=fillaccent2!45] (sw) at (0,3.4) {mjukvara på en processor};

  % Nåbarhetslagret.
  \node[blk, fill=fillaccent!45] (bridge) at (0,2.3) {seriebrygga};
  \node[blk, fill=fillaccent!45] (bank)   at (0,1.4) {registerbank};

  % CAN-sidan.
  \node[blk, fill=fillaccent3!50] (ctrl)  at (0,0.3) {kontroller\\\tiny tillståndsmaskin + arbitrering};
  \node[blk, minimum width=1.7cm] (timer) at (-2.1,-1.0) {bittimer};
  \node[blk, minimum width=1.7cm] (crc)   at (-0.1,-1.0) {CRC-15};
  \node[blk, minimum width=1.7cm] (shift) at (2.0,-1.0)  {skiftregister};

  \node[blk, fill=fillmuted] (bus) at (0,-2.2) {CAN-bussen};

  \draw[arr] (sw) -- (bridge);
  \draw[arr] (bridge) -- (bank);
  \draw[arr] (bank) -- (ctrl);
  \draw[arr] (ctrl) -- (timer);
  \draw[arr] (ctrl) -- (crc);
  \draw[arr] (ctrl) -- (shift);
  % Bussen drivs och läses av kontrollern, inte av lövblocken: routa runt dem.
  \draw[arr] (ctrl.east) -- ++(1.15,0) |- (bus.east);

  % Etiketter i marginalen.
  \node[anchor=west, muted, font=\scriptsize\itshape] at (2.9,1.85) {gör kontrollern nåbar};
  \node[anchor=west, muted, font=\scriptsize\itshape] at (2.9,-1.0) {känner till CAN};
\end{tikzpicture}
